\chapter{Padrões de Projeto e UML}
\label{cap:padroes-uml}

\begin{edital}
``Padrões de projeto'', ``design de software'', ``padrões de desenvolvimento
e reuso'', ``análise e projeto orientados a objetos''. UML não está listada
com todas as letras no edital do Perfil 3 (aparece explícita no Perfil 2), mas
é vizinha direta desses itens e a FGV cobra com frequência em provas de TI.
\end{edital}
\peso{ALTA PROBABILIDADE --- a FGV dá um cenário (``preciso garantir uma única
instância...'', ``preciso criar objetos sem acoplar à classe concreta...'') e
pergunta qual padrão resolve. Decore o gatilho de cada um, não só o nome.}

\section{Padrões de Projeto (Design Patterns / GoF)}

O enunciado descreve um \textbf{problema de design} e você identifica o
padrão que o resolve. A chave é reconhecer a \textbf{intenção} (o ``para
quê'') de cada padrão. O catálogo GoF (\emph{Gang of Four}) reúne
\textbf{23 padrões}, divididos em três grupos: \textbf{criacionais} (5),
\textbf{estruturais} (7) e \textbf{comportamentais} (11). O número e a
classificação em três grupos são cobrados diretamente.

Antes do catálogo, o essencial: um padrão de projeto \textbf{não é código
pronto, nem biblioteca, nem framework}. É a descrição de uma solução que já foi
reinventada tantas vezes que virou vocabulário --- e, como todo verbete, tem
quatro partes: um \textbf{nome}, o \textbf{problema} em que ele se aplica, a
\textbf{solução} em termos de classes e objetos, e as \textbf{consequências}
(porque toda escolha tem preço). Quem lê só a solução decora estrutura; quem lê
o problema reconhece o gatilho no enunciado.

Isso também explica o anti-padrão inverso: aplicar padrão \emph{sem} ter o
problema é \emph{over-engineering} --- uma fábrica abstrata para criar um único
tipo de objeto acrescenta duas classes e nenhuma flexibilidade.

\begin{conceito}[O que os 23 padrões têm em comum --- e o critério das três famílias]
Quase todo padrão do catálogo faz a \emph{mesma} coisa por dentro:
\textbf{isola o que varia atrás de uma interface estável} e prefere
\textbf{composição a herança} para poder trocar essa parte em tempo de
execução. Guarde essa frase: ela é o fio que liga Strategy, State, Bridge,
Decorator, Abstract Factory e quase todos os outros.

O que muda de família para família é \textbf{o que} está sendo isolado:

\begin{center}
\begin{tabular}{@{}p{3.2cm}p{3.8cm}p{4.8cm}@{}}
\toprule
\textbf{Família} & \textbf{Isola\ldots} & \textbf{Pergunta que ela responde} \\
\midrule
Criacionais (5) & \emph{qual} objeto se instancia, e quando & ``como criar isto sem amarrar o código à classe concreta?'' \\
Estruturais (7) & \emph{como} as peças se encaixam & ``como compor/embrulhar objetos para obter uma interface ou um comportamento?'' \\
Comportamentais (11) & \emph{quem} conversa com quem, e qual algoritmo roda & ``como distribuir a responsabilidade da interação em runtime?'' \\
\bottomrule
\end{tabular}
\end{center}

Na dúvida entre duas alternativas, classifique primeiro pela família: se o
enunciado fala em \textbf{criar} algo, os sete estruturais e os onze
comportamentais caem fora de uma vez.
\end{conceito}

\subsection{Criacionais (como criar objetos)}

\begin{center}
\begin{tabular}{@{}p{2.8cm}p{9cm}@{}}
\toprule
\textbf{Padrão} & \textbf{Intenção (gatilho no enunciado)} \\
\midrule
Singleton & garantir \textbf{uma única instância} e ponto global de acesso (``apenas uma conexão/configuração no sistema todo'') \\
Factory Method & delegar a criação a subclasses; criar objeto \textbf{sem acoplar à classe concreta} (``decidir em runtime qual objeto criar'') \\
Abstract Factory & criar \textbf{famílias} de objetos relacionados sem especificar classes concretas (``kit de UI para Windows e Mac'') \\
Builder & construir objeto complexo \textbf{passo a passo}, separando construção da representação (``montar um objeto com muitos parâmetros opcionais'') \\
Prototype & criar novos objetos \textbf{clonando} um protótipo (``copiar um objeto existente em vez de criar do zero'') \\
\bottomrule
\end{tabular}
\end{center}

\begin{pegadinha}
\textbf{Factory Method} (uma criação, via herança) $\times$ \textbf{Abstract
Factory} (famílias de produtos, via composição) --- a FGV troca os dois.
Singleton = \textbf{uma} instância, nunca ``mais de uma quando necessário''.
\end{pegadinha}

\subsection{Estruturais (como compor objetos/classes)}

\begin{center}
\begin{tabular}{@{}p{2.8cm}p{9cm}@{}}
\toprule
\textbf{Padrão} & \textbf{Intenção (gatilho)} \\
\midrule
Adapter & fazer interfaces \textbf{incompatíveis} trabalharem juntas (``adaptar uma API legada à nova'') \\
Facade & interface \textbf{simplificada} para um subsistema complexo (``uma fachada única para vários serviços'') \\
Proxy & um \textbf{substituto} que controla acesso ao objeto real (lazy load, cache, segurança) \\
Decorator & adicionar responsabilidades a um objeto \textbf{dinamicamente}, sem herança (``empilhar comportamentos em runtime'') \\
Composite & tratar objetos individuais e composições \textbf{de forma uniforme} (árvore parte-todo) \\
Bridge & separar \textbf{abstração} da \textbf{implementação} para variarem independentes \\
Flyweight & compartilhar objetos para economizar memória (muitos objetos semelhantes) \\
\bottomrule
\end{tabular}
\end{center}

\begin{pegadinha}
\textbf{Adapter} (compatibiliza o que já existe) $\times$ \textbf{Facade}
(simplifica um subsistema) $\times$ \textbf{Decorator} (adiciona
comportamento) --- os três se confundem na estrutura, mas a intenção é
diferente. \textbf{Proxy} controla \textbf{acesso}; \textbf{Decorator}
acrescenta \textbf{função}.
\end{pegadinha}

\subsection{Comportamentais (como objetos interagem)}

\begin{center}
\begin{tabular}{@{}p{2.8cm}p{9cm}@{}}
\toprule
\textbf{Padrão} & \textbf{Intenção (gatilho)} \\
\midrule
Strategy & família de algoritmos \textbf{intercambiáveis}, escolhidos em runtime (``trocar a regra de cálculo sem if gigante'') \\
Observer & notificar automaticamente \textbf{dependentes} quando o estado muda (publish-subscribe, eventos) \\
Command & encapsular uma \textbf{requisição como objeto} (undo/redo, fila de operações) \\
State & mudar o comportamento conforme o \textbf{estado interno} (``o objeto age diferente conforme sua situação'') \\
Template Method & definir o \textbf{esqueleto} de um algoritmo, deixando passos para subclasses \\
Iterator & percorrer uma coleção \textbf{sem expor sua estrutura} \\
Chain of Responsibility & passar a requisição por uma \textbf{cadeia} de handlers até alguém tratar \\
Mediator & centralizar a comunicação entre objetos (reduz acoplamento N-para-N) \\
Memento & capturar/restaurar o estado de um objeto (snapshot) \\
Visitor & adicionar operações a uma estrutura sem alterar as classes \\
Interpreter & definir uma \textbf{gramática} para uma linguagem e um interpretador que avalia suas sentenças (``interpretar expressões/regras escritas numa mini-linguagem'') \\
\bottomrule
\end{tabular}
\end{center}

\begin{pegadinha}
\textbf{Strategy} (troca o algoritmo) $\times$ \textbf{State} (troca conforme
o estado) --- parecidos na estrutura, diferentes na intenção.
\textbf{Observer} é o padrão de eventos/notificação; não confundir com
Mediator (que centraliza a comunicação, não notifica mudança de estado).
\textbf{Interpreter} é o 11º comportamental e o mais esquecido --- some da
lista de quem conta de cabeça, e é justamente aí que a FGV pergunta o número.
\end{pegadinha}

\textbf{Exemplo trabalhado --- do \texttt{if} gigante ao Strategy.} O cálculo
do frete começa assim, dentro do serviço de pedidos:

\begin{center}
\begin{tabular}{@{}p{5.6cm}p{6cm}@{}}
\toprule
\textbf{Antes} & \textbf{Depois} \\
\midrule
\texttt{if tipo == "sedex" \{\ldots\}} & interface \texttt{CalculoFrete} com um
método \texttt{calcular(pedido)} \\
\texttt{else if tipo == "pac" \{\ldots\}} & uma classe por regra:
\texttt{FreteSedex}, \texttt{FretePac}, \texttt{FreteRetirada} \\
\texttt{else if tipo == "retirada" \{\ldots\}} & o serviço \textbf{recebe} a
estratégia pronta e só chama \texttt{calcular} \\
\bottomrule
\end{tabular}
\end{center}

O que se ganhou tem nome: acrescentar a transportadora nova passou a ser
\textbf{criar uma classe}, não \textbf{editar} o método existente --- aberto à
extensão, fechado à modificação, que é o \emph{Open/Closed} do SOLID. E o
serviço deixou de conhecer as regras concretas: ele depende só da interface, o
que é o \emph{Dependency Inversion} do mesmo conjunto. Repare que o padrão não
eliminou a decisão de qual regra usar; ele a \textbf{moveu para fora}, para
quem monta o objeto.

Guarde o contraste, porque é a pegadinha do bloco: \textbf{Strategy e State têm
a mesma estrutura} --- um objeto delegando a outro por uma interface. O que
muda é quem decide a troca. No Strategy, \textbf{quem escolhe é o cliente}, de
fora, e as estratégias não sabem umas das outras. No State, \textbf{quem troca
é o próprio objeto} conforme seu estado interno, e os estados conhecem a
transição seguinte (``pedido pago passa a enviado''). Se o enunciado descreve
um ciclo de vida com transições, é State; se descreve alternativas
intercambiáveis para a mesma tarefa, é Strategy.

\subsection{GRASP}

O edital cita GRASP no Perfil 2, mas o conceito permeia o Perfil 3. São
princípios (não ``caixas de código'') para decidir \textbf{qual classe recebe
qual responsabilidade}: Information Expert, Creator, Controller, Low
Coupling, High Cohesion, Polymorphism, Pure Fabrication, Indirection,
Protected Variations.

Dois deles respondem sozinhos quase toda questão de ``a qual classe cabe esta
responsabilidade''. \textbf{Information Expert:} a responsabilidade vai para a
classe que \textbf{tem os dados} necessários para cumpri-la --- quem calcula o
total do pedido é o \texttt{Pedido}, que conhece os itens, não um
\texttt{CalculadoraDeTotais} que precisaria pedir tudo emprestado.
\textbf{Creator:} quem cria o objeto é quem o \textbf{agrega, contém ou usa de
perto} --- o \texttt{Pedido} cria o \texttt{ItemDePedido}. Os dois existem pelo
mesmo motivo: manter juntos o dado e o comportamento que o usa é o que produz
\textbf{alta coesão} e \textbf{baixo acoplamento}, que são eles próprios
princípios da lista.

\subsection{Padrões arquiteturais (não são GoF, mas caem)}

\begin{itemize}
\item \textbf{MVC} (Model-View-Controller): separa dados, apresentação e
controle.
\item \textbf{MVVM / MVP}: variações com binding/presenter.
\item \textbf{DAO / Repository}: isola o acesso a dados.
\item \textbf{DTO}: objeto para transportar dados entre camadas.
\item \textbf{Front Controller}, Singleton de configuração --- aparecem em
Java EE.
\end{itemize}

\subsection{Reuso e anti-padrões}

\textbf{Reuso:} herança (``é-um'') $\times$ \textbf{composição} (``tem-um'')
--- prefira composição (\emph{favor composition over inheritance}). Também:
bibliotecas, frameworks, componentes.

\textbf{Anti-padrões:} God Object (classe que faz tudo), Spaghetti Code,
Golden Hammer (usar sempre a mesma solução), Copy-Paste. A FGV pergunta ``qual
é o anti-padrão'' ou ``qual prática viola SOLID''.

\begin{conceito}[SOLID --- os cinco, para consultar aqui]
Os princípios aparecem \textbf{com todas as letras no edital do Perfil 2}
(item 16) e são a linguagem em que a banca cobra ``qual prática viola'' e
``que princípio o padrão X concretiza''. O tratamento completo, com as
armadilhas de Java, está no Capítulo \ref{cap:java}; aqui fica o essencial
para ligar padrão a princípio sem sair do capítulo:

\begin{itemize}
\item \textbf{S} --- \emph{Single Responsibility}: uma responsabilidade (um
motivo para mudar) por classe.
\item \textbf{O} --- \emph{Open/Closed}: aberta à \textbf{extensão}, fechada à
\textbf{modificação}.
\item \textbf{L} --- \emph{Liskov}: o subtipo substitui o tipo base sem
quebrar o programa.
\item \textbf{I} --- \emph{Interface Segregation}: interfaces pequenas e
específicas, em vez de uma interface gorda.
\item \textbf{D} --- \emph{Dependency Inversion}: depender de
\textbf{abstração}, não de implementação concreta.
\end{itemize}

O par que a FGV inverte é \textbf{ISP} (interface enxuta) $\times$
\textbf{SRP} (uma responsabilidade) --- os dois falam em ``pequeno e
específico'', mas um trata de \emph{interface} e o outro de \emph{classe}.
\end{conceito}

\begin{comosair}
\begin{enumerate}
\item Leia a \textbf{intenção}, não a implementação. O enunciado descreve o
problema; case com o ``para quê'' das tabelas acima.
\item Memorize os pares confundíveis: Factory Method $\times$ Abstract
Factory; Adapter $\times$ Facade $\times$ Decorator; Strategy $\times$ State;
Proxy $\times$ Decorator.
\item Ligue ao SOLID: Strategy e Template Method concretizam o Open/Closed;
Dependency Inversion aparece em Factory/Abstract Factory.
\end{enumerate}
Os mais cobrados pela FGV: \textbf{Singleton, Factory Method/Abstract
Factory, Strategy, Observer, Adapter, Facade, Decorator, MVC.} O Spring usa
Singleton (beans), Factory, Proxy (AOP), Template Method (JdbcTemplate),
Observer (eventos) --- ligar padrão a framework ajuda a fixar.
\end{comosair}

\begin{jacaiu}
\textbf{Em prova real da FGV:} \textbf{uma} questão na amostra inteira ---
\textbf{ALERO 2026}, cache de documentos que precisa de uma \emph{única
instância ativa} em toda a aplicação, gabarito \textbf{Singleton}, com Factory
Method, Builder, Observer e Strategy como distratores. O enunciado fala em
``Padrão de Projeto da \textbf{classificação GoF}'' --- por isso a tabela das
três famílias acima continua valendo o estudo: é assim que a banca ancora o
item.

\textbf{No nosso banco} (previsto pelo edital, ainda não visto na amostra de
provas --- 23 questões, subtag \texttt{padroes-projeto}): os demais
criacionais em cenário (Abstract Factory para famílias de componentes; Builder
com muitos atributos opcionais; Prototype quando criar do zero é caro; Factory
Method); Adapter para biblioteca de terceiros com interface incompatível;
Proxy como substituto de imagem pesada; Decorator para responsabilidade
dinâmica; Composite em árvore de pastas e arquivos; Facade acionando vários
subsistemas em sequência; Strategy trocando o cálculo de frete em runtime;
Observer na planilha que recalcula dependentes; Template Method nas etapas
compartilhadas; Command na fila com desfazer; Chain of Responsibility nas
validações sucessivas; State no pedido que muda de comportamento por estágio;
MVC; e GRASP (a qual classe cabe a responsabilidade).

\textbf{Cuidado ao garimpar ``já caiu'' por palavra-chave:} o corpus tem
quatro falsos positivos que \emph{não} são GoF --- ``Decorators (@)'' numa
questão de Python (é metaprogramação da linguagem), ``Portas e Adaptadores''
na arquitetura hexagonal (não é o Adapter), os vários ``proxy'' das questões
de rede, e o MVC, que só aparece descrito sem ser nomeado. Rode
\texttt{./quiz.py padroes-projeto}.
\end{jacaiu}

\section{UML}
\label{sec:uml}

A UML é uma \textbf{linguagem de notação, não um método}: ela padroniza o
desenho, não diz quando desenhar --- isso é papel do processo (RUP, Scrum;
Capítulo \ref{cap:eng-software}). O problema que ela resolve é banal e caro:
antes dela, cada empresa tinha a própria simbologia, e o mesmo losango
significava coisas diferentes em dois documentos da mesma sala. Padronizar a
notação é o que permite discutir projeto sem discutir desenho. A especificação
vigente da OMG é a \textbf{2.5.1}, e a FGV costuma imprimir esse número no
enunciado --- é sinal de que o item vai cobrar a semântica exata de um símbolo.

\subsection{Os 14 diagramas em 2 famílias}

\begin{center}
\begin{tabular}{@{}p{2.6cm}p{2.6cm}p{6.4cm}@{}}
\toprule
\textbf{Família} & \textbf{Foca em} & \textbf{Diagramas principais} \\
\midrule
Estruturais & estrutura estática & \textbf{Classe}, Objeto, Componente, Implantação (deployment), Pacote, Estrutura composta, Perfil \\
Comportamentais & comportamento dinâmico & \textbf{Caso de uso}, \textbf{Sequência}, \textbf{Atividade}, \textbf{Estado} (máquina de estados), Comunicação, Interação geral, Tempo \\
\bottomrule
\end{tabular}
\end{center}

\begin{pegadinha}
A FGV troca estrutural $\leftrightarrow$ comportamental. \textbf{Classe} é
estrutural; \textbf{sequência, atividade, estado, caso de uso} são
comportamentais. Nunca o contrário.
\end{pegadinha}

\subsection{Diagrama de Classes (o mais cobrado)}

Mostra classes, atributos, métodos e relacionamentos:

\begin{center}
\begin{tabular}{@{}p{3.6cm}p{4.6cm}p{2.6cm}@{}}
\toprule
\textbf{Relacionamento} & \textbf{Significado} & \textbf{Notação} \\
\midrule
Associação & ligação entre classes & linha \\
Agregação & todo-parte \textbf{fraca} (a parte vive sem o todo) & losango \textbf{vazio} \\
Composição & todo-parte \textbf{forte} (a parte morre com o todo) & losango \textbf{cheio} \\
Herança / Generalização & é-um (subtipo) & seta triângulo vazio \\
Dependência & usa temporariamente & seta tracejada \\
Realização & implementa interface & tracejada + triângulo \\
\bottomrule
\end{tabular}
\end{center}

\textbf{Multiplicidade:} \texttt{1}, \texttt{0..1}, \texttt{1..*}, \texttt{*}
(muitos). \textbf{Visibilidade:} \texttt{+} público, \texttt{-} privado,
\texttt{\#} protegido, \texttt{\~{}} pacote.

\begin{conceito}[Ler a caixa e a linha --- a sintaxe que a FGV cobra de perto]
A classe tem \textbf{três compartimentos}: nome, atributos, operações. E o
atributo tem uma forma completa, que é onde mora o item de leitura:

\begin{center}
\texttt{- status : String [1] = "aberto"}
\end{center}

\noindent visibilidade, nome, tipo, multiplicidade e \textbf{valor padrão} ---
o valor que a instância assume quando ninguém informa outro na criação.

Do lado da linha, duas notações que aparecem como distrator justamente por
serem pouco lidas:

\begin{itemize}
\item \textbf{Extremidade de associação $\times$ atributo.} Cada ponta da
associação tem papel, multiplicidade e navegabilidade --- e uma ponta navegável
\textbf{pode ser representada como atributo} da classe do outro lado. ``Pedido
associa-se a 1 Cliente'' e ``Pedido tem o atributo \texttt{cliente: Cliente}''
dizem a \emph{mesma} coisa em UML; são duas notações para um modelo só, não
duas modelagens diferentes.
\item \textbf{Associação qualificada.} Um qualificador (uma chave, desenhado
num retângulo na ponta) \textbf{particiona} o conjunto do outro lado e
tipicamente derruba a multiplicidade de \texttt{*} para \texttt{0..1}: dado um
banco \emph{e} um número de conta, chega-se a no máximo uma conta. É um
índice, não um todo-parte --- e é por isso que ela funciona como distrator de
agregação.
\end{itemize}

\textbf{Vocabulário da especificação:} o que este capítulo chama de composição
a UML chama de \textbf{agregação composta} (o losango cheio). Se a alternativa
usar esse nome, é composição --- não é um terceiro tipo de relacionamento.
\end{conceito}

\begin{pegadinha}
Pegadinha central do bloco: \textbf{agregação} (losango vazio, parte
independente) $\times$ \textbf{composição} (losango cheio, parte dependente).
A FGV inverte qual dos dois ``morre com o todo''.
\end{pegadinha}

\subsection{Casos de Uso (requisitos funcionais)}

Ator (boneco), caso de uso (elipse), fronteira do sistema.

\begin{itemize}
\item \textbf{«include»}: comportamento \textbf{sempre} incluído
(obrigatório).
\item \textbf{«extend»}: comportamento \textbf{opcional/condicional}.
\end{itemize}

\begin{pegadinha}
\textbf{include} (sempre executa) $\times$ \textbf{extend} (às vezes executa)
--- a FGV inverte esse par com frequência.
\end{pegadinha}

\subsection{Sequência (interação no tempo)}

Linhas de vida (objetos no topo), mensagens trocadas na ordem temporal (de
cima para baixo), barras de ativação. Mensagem síncrona (seta cheia) $\times$
assíncrona (seta aberta) $\times$ retorno (tracejada).

\begin{conceito}[Os quatro diagramas de interação --- onde a FGV monta os distratores]
Quatro diagramas mostram \textbf{objetos trocando mensagens}. Como todos
``mostram interação'', a banca usa os outros três como distratores do de
sequência --- e a diferença não é o conteúdo, é \textbf{o que cada um
enfatiza}:

\begin{center}
\begin{tabular}{@{}p{3.2cm}p{4.2cm}p{4.2cm}@{}}
\toprule
\textbf{Diagrama} & \textbf{Enfatiza} & \textbf{Como se lê a ordem} \\
\midrule
\textbf{Sequência} & a \textbf{ordem cronológica} das mensagens & pela posição
vertical: de cima para baixo \\
\textbf{Comunicação} & a \textbf{estrutura de ligações} entre os objetos (quem
está conectado a quem) & pela \textbf{numeração} das mensagens (1, 1.1,
1.2) --- não há eixo de tempo \\
\textbf{Tempo} (\emph{timing}) & a \textbf{mudança de estado} de cada
participante ao longo de uma \textbf{escala de tempo explícita} & pelo eixo
horizontal de tempo; usado quando o requisito é de duração/tempo real \\
\textbf{Visão geral de interação} & como \textbf{vários cenários} se encadeiam
& é um diagrama de atividades cujos nós são interações \\
\bottomrule
\end{tabular}
\end{center}

A regra de decisão cabe numa linha: pediu a \textbf{ordem exata em que as
mensagens acontecem}, é \textbf{sequência}. Sequência e comunicação são
\textbf{semanticamente equivalentes} (dá para converter um no outro sem perder
informação) --- o que muda é a ênfase, e é justamente aí que a alternativa
``comunicação'' parece defensável e não é, quando o enunciado fala em
cronologia.

Não confunda com \textbf{atividade}: lá não há troca de mensagens entre
participantes, há \textbf{fluxo de trabalho} --- ações, decisões e paralelismo.
\end{conceito}

\subsection{Atividade e Estado}

\textbf{Atividade:} fluxo de trabalho/algoritmo --- nós de ação, decisão
(losango), bifurcação/junção (barra, paralelismo), início/fim. Parecido com
fluxograma; usado para modelar processos.

\textbf{Máquina de estados:} os estados de um objeto e as transições
disparadas por eventos (o objeto reage diferente conforme o estado --- casa
com o padrão de projeto State).

\begin{comosair}
\begin{enumerate}
\item Decore a família de cada diagrama (estrutural $\times$
comportamental) --- é a pegadinha mais comum.
\item No diagrama de classes, memorize agregação $\times$ composição (vazio
$\times$ cheio) e as multiplicidades.
\item Caso de uso: include (obrigatório) $\times$ extend (opcional).
\item Ligue ao resto: caso de uso $\leftrightarrow$ requisitos funcionais;
classes $\leftrightarrow$ OO/SOLID; estado $\leftrightarrow$ padrão State.
\end{enumerate}
Diagramas mais cobrados: \textbf{Classe, Caso de Uso, Sequência, Atividade.}
Deployment e Componente aparecem em questões de arquitetura (Capítulo
\ref{cap:arquitetura}).
\end{comosair}

\begin{jacaiu}
\textbf{Em prova real da FGV:} ao contrário dos padrões, UML \textbf{cai}, e
sempre citando a versão. \textbf{Diagrama de sequência} para a ordem
cronológica exata das mensagens --- e os distratores eram os outros diagramas
de interação (\textbf{comunicação}, \textbf{timing}, \textbf{visão geral de
interação}) mais o de atividades; \textbf{composição} para a parte que não
existe sem o todo e some junto com ele (distratores: agregação, dependência,
associação qualificada); \textbf{diagrama de casos de uso} combinado com o
\textbf{diagrama de objetos} para montar plano de teste (o de objetos mostra
os \emph{valores} dos atributos das instâncias num ponto da execução) ---
\textbf{ALERO 2026}, sempre com ``UML 2.5.1'' no enunciado. Leitura de
\textbf{diagrama de classes} (extremidade de associação representável como
atributo, agregação composta, valor default de atributo) --- \textbf{MPU}.
\textbf{Generalização/especialização total e disjunta} caiu na
\textbf{ALERO 2026} pelo lado do modelo ER, com a mesma semântica.

\textbf{No nosso banco} (previsto pelo edital, ainda não visto na amostra de
provas --- 10 questões, subtag \texttt{uml}): as duas famílias na UML 2.x;
multiplicidade em diagrama de classes; \texttt{<<include>>} $\times$
\texttt{<<extend>>}; atividade para decisão, paralelismo e sincronização;
máquina de estados no ciclo de vida do objeto; direção da seta da
generalização; componentes $\times$ implantação; realização de interface.

Rode \texttt{./quiz.py uml}.
\end{jacaiu}
